\section{Preliminary Evaluation}\label{sec:eval}
\textbf{RQ1 (H0):} how often do LLMs misexplain an intervention, ungrounded and grounded?
\textbf{RQ2 (H1, model-side):} are counterfactual questions harder than matched descriptive ones?

\textbf{Corpus.} We sampled 88 HumanEval+ and 88 MBPP+ tasks~\cite{chen2021humaneval,austin2021mbpp,liu2023evalplus}
(seeded), with canonical solutions as programs and up to 120 EvalPlus inputs as $T$ (mean
\SMeanInputsPerFunction). Of \SNFunctionsAttempted{} functions, \SNNoCandidates{} had no mutation site, \SNNoNonequivalent{}
had no non-equivalent candidate among those tried, and \SNMutants{} received a mutant (Table~\ref{tab:corpus}; error columns descriptive, G$^{*}$:
copies count as errors);
\PEquivalentCandidateRate\% of tried candidates were equivalent on $T$ and \SNDegenerate{} degenerate. Most mutants are
broad (median \PzMedianDifferingFraction\% of inputs differ; \SNMutantsAllDiffer{} differ on all); only removed checks are
typically narrow; \SNMutantsRaiseMostly{} raise an exception on most inputs.
\begin{table}[t]
\centering\footnotesize
\caption{Corpus by operator ($n$: HE+/MBPP+). diff: median share of inputs whose result changes; rais.: raising on most inputs; witness/affected: Qwen2.5-Coder-14B error (\%), U/G$^{*}$, U/G.}
\label{tab:corpus}
\setlength{\tabcolsep}{3pt}
\begin{tabular}{@{}lrcrcc@{}}
\toprule
Operator & $n$ & diff & rais. & witness & affected \\
\midrule
negate guard & 19 (16/3) & 0.95 & 2 & 63/79 & 100/95 \\
remove check & 13 (11/2) & 0.06 & 0 & 77/77 & 31/23 \\
comparison op & 25 (11/14) & 0.97 & 1 & 64/68 & 92/88 \\
arithmetic op & 52 (26/26) & 0.95 & 3 & 81/69 & 77/67 \\
boundary const & 20 (12/8) & 0.67 & 0 & 60/65 & 85/65 \\
\midrule
\textbf{all} & 129 (76/53) & 0.91 & 6 & 71/71 & 80/71 \\
\bottomrule
\end{tabular}
\end{table}


\textbf{Generators and statistics.} \GenModel{} (local, temperature 0, fixed seed) explains all functions;
\StrongModel{} (\texttt{claude-haiku-5-5}, default sampling) explains all functions and writes predicates for a seeded
random subset of \SPredStrongN{} (API budget: \SLlmCallsStrong{} calls in total). All \SLlmCacheFiles{} responses are
cached; the analysis replays exactly. CIs are percentile bootstraps over functions; tests are exact
McNemar~\cite{mcnemar1947note} with Holm correction over a family of \SHolmFamilySize{} tests fixed in a second analysis
phase (witness raw and copies-as-error, affected-set exact, per generator; descriptive and $P_m$-shown vs.\ counterfactual,
per probe model; the docstring ablation). In Table~\ref{tab:h0}, affected-set error = $\geq$1 misclassified input
(unscorable = error); bal.\ acc.\ (unscorable = 0.5) and F1 (scorable only) are per-function means; const.\ = ``every
input is affected'' (all functions).
\begin{table}[t]
\centering\footnotesize
\caption{H0: ungrounded (U) vs.\ grounded (G), paired by function, in \%; $\Delta$=U$-$G, 95\% bootstrap CI; $p$: exact McNemar, Holm-adjusted ($^{\dagger}$: unadjusted, secondary; bal.\ acc.: Wilcoxon).}
\label{tab:h0}
\setlength{\tabcolsep}{3pt}
\begin{tabular}{@{}lrrcr@{}}
\toprule
Measure & U & G & $\Delta$ & $p$ \\
\midrule
\multicolumn{5}{@{}l}{\emph{Qwen2.5-Coder-14B} (witness $n$=129; affected $n$=129)} \\
\quad witness error & 71.3 & 44.2 & 27.1 [16, 37] & $<$0.001 \\
\quad witness, copies = error & 71.3 & 70.5 & 0.8 [-9, 11] & 1.000 \\
\quad \ \ malformed not counted & 60.5 & 58.9 & 1.6 [-9, 11] & 0.880$^{\dagger}$ \\
\quad affected-set error & 79.8 & 70.5 & 9.3 [2, 16] & 0.118 \\
\quad \ \ scorable only (3/5 unscor.) & 79.4 & 69.4 & & \\
\quad affected bal.\ acc.\ (const.\ 59.7) & 65.8 & 72.0 & -6.2 & 0.110$^{\dagger}$ \\
\quad affected F1 (const.\ 73.6) & 58.6 & 69.9 & & \\
\multicolumn{5}{@{}l}{\emph{Claude Haiku 5.5} (witness $n$=129; affected $n$=72)} \\
\quad witness error & 3.9 & 2.3 & 1.6 [-2, 5] & 1.000 \\
\quad witness, copies = error & 3.9 & 3.1 & 0.8 [-2, 4] & 1.000 \\
\quad \ \ malformed not counted & 3.1 & 3.1 & 0.0 [-3, 3] & 1.000$^{\dagger}$ \\
\quad affected-set error & 5.6 & 1.4 & 4.2 [0, 10] & 1.000 \\
\quad \ \ scorable only (2/0 unscor.) & 2.9 & 1.4 & & \\
\quad affected bal.\ acc.\ (const.\ 61.1) & 97.4 & 99.4 & -2.0 & 0.066$^{\dagger}$ \\
\quad affected F1 (const.\ 76.5) & 98.6 & 99.5 & & \\
\bottomrule
\end{tabular}
\end{table}


\textbf{RQ1: grounding helps only by copying.} For \GenModel{}, \PzWitGenRawA\% of ungrounded witnesses fail the check
(\PzFailGenUngroundedMalformedwitness\% malformed argument lists); grounded, \PzWitGenRawB\%. But \PzWitGenCopied\% of grounded witnesses were copied from the evidence despite the
instruction (\SWitGenCopiedPass{} of the grounded passes); the model's own valid witnesses go from \SWitGenOwnValidU{}
to \SWitGenOwnValidG{}, and counting copies as errors gives \PWitGenStrictB\% vs.\ \PWitGenStrictA\% ($p$=\SWitGenStrictPHolm,
Table~\ref{tab:h0}). \GenModel's affected-input predicates are poor: on exact match they misclassify some test input for
\PzPredGenExactA\% of functions ungrounded, close to the constant ``every input is affected'' predicate
(\PzPredGenTrivialExactErr\%); on F1 they are worse than that constant (\PPredGenUFone{} vs.\ \PPredGenTrivialFone), and only
balanced accuracy puts them above it (\PPredGenBalaccU{} vs.\ \PPredGenTrivialBal). Grounded, exact-match error is
\PzPredGenExactB\% ($p$=\SPredGenExactP, Holm \SPredGenExactPHolm), but the evidence block also reveals $|\Delta_T|/|T|$:
for functions that differ on every input it says so, and four grounded predicates are literally \texttt{return True};
without those functions the contrast is \SPredGenExactNotalldifferNonezero/\SPredGenExactNotalldifferNzeroone{}
($p$=\SPredGenExactNotalldifferP), the predicate analogue of copying. As a non-LLM witness baseline, seeded type-preserving differential fuzzing (50 fresh inputs per mutant) found a valid witness on 102/129 (79.1\%) mutants within 10 inputs, 110/129 (85.3\%) within 25, and 113/129 (87.6\%) within 50. On up to 50 fresh inputs strictly outside $T$ per function, of predicates exact on $T$, \StrongModel{} remained exact on all fresh inputs for 65/68 (U) and 64/71 (G) functions, whereas \GenModel{} remained exact for 17/26 (U) and 21/38 (G).
\begin{figure}[t]
\centering
\begin{tikzpicture}
\begin{axis}[width=\columnwidth,height=4.6cm,ybar=1pt,bar width=7pt,ymin=0,ymax=118,
  ylabel={affected-set error (\%)},ylabel style={font=\footnotesize},
  symbolic x coords={narrow,mid,broad},xtick=data,xticklabels={{$<$10\% ($n$=19)},{10--50\% ($n$=24)},{$\geq$50\% ($n$=86)}},
  x tick label style={font=\scriptsize},y tick label style={font=\scriptsize},enlarge x limits=0.22,
  nodes near coords,nodes near coords style={font=\tiny,rotate=90,anchor=west,/pgf/number format/.cd,fixed,precision=1,fixed zerofill},
  legend style={at={(0.5,1.02)},anchor=south,legend columns=2,font=\scriptsize,draw=none,/tikz/every even column/.append style={column sep=4pt}},
  legend image code/.code={\draw[#1] (0cm,-0.08cm) rectangle (0.3cm,0.12cm);},
  ymajorgrids,grid style={dashed,gray!30}]
\addplot[ybar,fill=white,draw=black] coordinates {(narrow,68.4) (mid,91.7) (broad,79.1)};
\addlegendentry{Qwen2.5-Coder-14B, U}
\addplot[ybar,fill=gray!55,draw=black] coordinates {(narrow,63.2) (mid,83.3) (broad,68.6)};
\addlegendentry{Qwen2.5-Coder-14B, G}
\addplot[ybar,fill=white,draw=black,postaction={pattern=north east lines}] coordinates {(narrow,0.0) (mid,15.4) (broad,4.0)};
\addlegendentry{Claude Haiku 5.5, U}
\addplot[ybar,fill=gray!55,draw=black,postaction={pattern=crosshatch}] coordinates {(narrow,0.0) (mid,7.7) (broad,0.0)};
\addlegendentry{Claude Haiku 5.5, G}
\end{axis}
\end{tikzpicture}
\caption{Affected-set error by breadth of the executed difference (fraction of test inputs whose result changes); tick labels give Qwen2.5-Coder-14B's $n$. Claude Haiku 5.5 bars cover its 72-function subset ($n$=9/13/50). Values printed on bars (0.0 = no error).}
\label{fig:strata}
\end{figure}


\textbf{LLM judges} (\JudgeA, \JudgeB) graded the $A$/$d$ prose and are not evidence: \JudgeA{} rejects
\PzJoneRejectRateGen\% of \GenModel{} and \PzJoneRejectRateStrong\% of \StrongModel{} explanations; \JudgeB{} is unvalidated ($\kappa$=\SJtwoProgKappaGen{}
with the witness check).

\textbf{RQ2 pilot (models as stand-ins, NOT human data).} A model answers, on the shortest differing input $x^\star$, a \emph{descriptive}
(what does $P(x^\star)$ return?) and a \emph{counterfactual} question (given the diff, what would the modified function
return?), the counterfactual question on an \emph{unaffected} input, and the descriptive question with $P_m$ shown. An initial run kept HumanEval+ docstrings (which
describe the \emph{original}) and cited raw-source line numbers, wrong for \SNLineMismatchRoundone{} prompts; the
main run strips docstrings and comments, renames the entry function \texttt{f} (four helpers keep their names) and
cites lines in the shown code (behaviour verified unchanged by execution). The ablation in Table~\ref{tab:probe} confirms the confound: on HumanEval+, accuracy with $P_m$
shown rises from \PzAblPonefourDESCMUTHeB\% to \PzAblPonefourDESCMUTHeA\% once docstring, comments and name are removed ($p$=\SAblPonefourDESCMUTHePHolm),
and answers that \emph{revert} to the original's output~\cite{yang2025reversion} fall from \PzProbePonefourdDESCMUTHeReverts\%
to \PzProbePonefoursDESCMUTHeReverts\%, consistent with CodeCrash~\cite{lam2025codecrash}; MBPP+ (no docstrings) barely moves.
With cues removed, \GenModel{} answers \PzProbePonefoursDESC\% of descriptive and \PzProbePonefoursCF\% of counterfactual
questions ($\Delta$=\PProbePonefoursDESCVsCFD{} points, $p$\,\SProbePonefoursDESCVsCFPHolm). On unaffected inputs it is
right \PzProbePonefoursCFCTRL\% of the time, close to its descriptive accuracy on the same functions
(\PzProbePonefoursDESCOnCtrl\%), so the diff presentation itself is not the difficulty; this control alone cannot rule out
ignoring the edit. What argues against pure inertia is that only \PzProbePonefoursCFReverts\% of counterfactual answers on
affected inputs revert to the original's output (inertia predicts $\approx$\PzProbePonefoursDESC\%): treating ``the answer
changes'' as a detection, the hit rate is \PProbePonefoursHit\% vs.\ a false-alarm rate of \PProbePonefoursFa\%
($d'$=\SProbePonefoursDprime; hits over 129, false alarms over 104 functions; \ProbeB{} \PProbePsevensHit\% vs.\ \PProbePsevensFa\%, $d'$=\SProbePsevensDprime, closer to
inertia). With $P_m$ shown in full, accuracy (\PzProbePonefoursDESCMUT\%) is not detectably different from the
counterfactual question ($p$=\SProbePonefoursDESCMUTVsCFPHolm): most of the difficulty is tracing mutated code, consistent
with memorised originals~\cite{yang2025reversion}, not the counterfactual framing. \ProbeB{} shows the same ordering
(\PzProbePsevensDESC\% vs.\ \PzProbePsevensCF\%, $p$\,\SProbePsevensDESCVsCFPHolm). Humans have not memorised the originals, so this is a weak prior for H1.
\begin{table}[t]
\centering\footnotesize
\caption{Model-side probe (NOT human data): accuracy (\%) by benchmark; main rows without docstrings, comments and entry name. HE+ $n$=76, MBPP+ $n$=53 (CF on an unaffected input: $n$=69/35).}
\label{tab:probe}
\setlength{\tabcolsep}{3pt}
\begin{tabular}{@{}lcccc@{}}
\toprule
 & \multicolumn{2}{c}{Qwen2.5-Coder-14B} & \multicolumn{2}{c}{Qwen2.5-Coder-7B} \\
Question & HE+ & MBPP+ & HE+ & MBPP+ \\
\midrule
descriptive ($P$) & 79 & 75 & 78 & 74 \\
counterfactual & 46 & 51 & 33 & 34 \\
counterf., unaffected $x$ & 75 & 66 & 75 & 63 \\
descriptive ($P_m$ shown) & 46 & 62 & 36 & 51 \\
\midrule
\multicolumn{5}{@{}l}{\emph{Ablation: initial materials (docstring, comments and name kept)}} \\
descriptive ($P$) & 83 & 79 & 76 & 75 \\
descriptive ($P_m$ shown) & 29 & 58 & -- & -- \\
\bottomrule
\end{tabular}
\end{table}

